% !TEX root = ../main.tex
% =========================================================
% abstrak.tex
% Abstrak Laporan Kerja Praktik (Indonesia + Inggris)
% =========================================================

% ---------------- ABSTRACT (Inggris) ----------------
\clearpage
\phantomsection
\addcontentsline{toc}{chapter}{ABSTRACT}
\begin{center}
    \textbf{ABSTRACT}\\[0.5cm]
\end{center}

Floating photovoltaic (PV) power plants operate in environments influenced by
water and humidity, so PV modules may experience abnormal conditions,
including submersion. A module that has been submerged cannot be declared
serviceable based on visual inspection alone; an electrical-parameter based
evaluation is therefore required. This internship report analyses the
post-submersion serviceability of PV modules at the Cirata Floating PV plant
using current--voltage (I--V) curve testing and current--voltage indicators.
The analysis covers 59 test records of JKM560N-72HL4-BDV modules measured
onshore with an I--V curve tracer, with parameters corrected to Standard Test
Conditions (STC). The evaluation compares STC parameters against the nameplate,
computes the Fill Factor, and analyses the $R_{VM}$ and $R_{IM}$ indicators to
draw an initial interpretation for operation and maintenance activities. The
results show that 6 records (10.17\%) obtained an Ok status while 53 records
(89.83\%) obtained a Not Ok status, with the Not Ok group exhibiting an average
STC maximum power of about 501.6 W, approximately 10.4\% below the 560 W
nameplate value. The $R_{VM}$ indicator separates the two groups more clearly
than $R_{IM}$, indicating that the performance difference is dominated by the
voltage side. Because part of the Ok data was obtained under low irradiance,
the serviceability status must be interpreted carefully and confirmed through
retesting under more representative irradiance as well as an insulation
resistance test before any reinstatement decision is made.

\vspace{1em}
\noindent Key words: floating PV plant, submerged PV module, I--V testing,
current--voltage indicators, Fill Factor

